\bookchapter{Capítulo 6: Protege tu vida digital — Seguridad y privacidad}

\chapterintro{Seguridad del teléfono = seguridad de la privacidad + seguridad patrimonial. En esta era en que los datos son un activo, proteger el teléfono es protegerte a ti mismo.}

\booksection{I. Bloqueo de pantalla: la primera línea de defensa física}

El bloqueo de pantalla es la primera línea de defensa de tu teléfono y también la más fácil de pasar por alto.

\booksubsection{Clasificación por nivel de seguridad:}

\begin{enumerate}
  \item \textbf{Contraseña/PIN}: máxima seguridad, es el respaldo final (cuando fallan la huella o el rostro).
  \item \textbf{Huella dactilar}: equilibrio perfecto entre seguridad y comodidad.
  \item \textbf{Reconocimiento facial}: lo más cómodo, pero la seguridad depende de la tecnología (3D estructurado es seguro; 2D con foto se engaña fácilmente).
  \item \textbf{Patrón de desbloqueo}: cómodo, pero deja huellas que pueden adivinarse.
\end{enumerate}

\booksubsection{Configura un bloqueo fuerte:}

\begin{itemize}
  \item Establece una contraseña compleja (no uses tu cumpleaños ni 123456).
  \item Activa la biometría.
  \item Ajusta el tiempo de bloqueo a "Inmediato" o al más corto posible.
\end{itemize}

\booksection{II. “Buscar mi teléfono”: la herramienta antirrobo definitiva}

\textbf{¡Debe estar activada!}Permite:

\begin{itemize}
  \item \textbf{Localizar}: ver en el mapa dónde se perdió el teléfono.
  \item \textbf{Reproducir sonido}: hacer sonar el teléfono, útil para encontrarlo en casa.
  \item \textbf{Modo perdido}: bloquear el teléfono de forma remota y mostrar un número de contacto.
  \item \textbf{Borrado remoto}: último recurso, borra todos los datos (tras el borrado ya no se puede localizar).
\end{itemize}

En [Configuración], Apple busca “Buscar mi iPhone”; Android, “Buscar mi dispositivo” o un servicio similar de la marca; asegúrate de que esté activado (en verde).

\booksection{III. Actualizaciones del sistema: parchear y tapar huecos}

La finalidad principal de las actualizaciones del sistema es reparar vulnerabilidades de seguridad, como tapar grietas en una casa. Por eso, al ver un aviso de actualización, conviene actualizar cuanto antes.

Eso sí, las grandes actualizaciones de versión (por ejemplo, de iOS 17 a iOS 18) pueden tener problemas de compatibilidad; puedes esperar unos días antes de actualizar.

\booksection{IV. Identificar y prevenir el fraude telefónico}

En una era de rápido desarrollo de la IA, el modelo tradicional de “ver para creer” ha perdido vigencia. Los estafadores usan tecnologías de IA como el deepfake para pasar del fraude “de escopeta” al “acecho selectivo” y la depredación precisa.

\booksubsection{Los timos más habituales en la era de la IA}

\textbf{(1) Fraude con intercambio de rostro y voz por IA}\par

Los delincuentes recopilan fotos en alta resolución o muestras de voz de pocos segundos publicadas por el objetivo en redes sociales y, con IA de deepfake, clonan un “video de un ser querido” o una “voz” suficientemente creíbles. Guiones habituales:\textbf{accidente repentino, petición urgente de dinero, secuestro}, etc. Explotan los lazos afectivos y familiares, por lo que\textbf{son extremadamente engañosos}.

\textbf{(2) Suplantación de autoridades / servicio al cliente}\par

La IA puede simular con precisión la voz y el tono de agentes de servicio al cliente oficiales o de autoridades policíales y judiciales, e incluso sintetizar capturas de transferencia, documentos oficiales y órdenes de detención altamente realistas,\textbf{generando pánico o presión de autoridad}para inducir a la víctima a transferir dinero o activar la pantalla compartida y robar así la información de la cuenta.

\textbf{(3) Fraude de inversión con IA}\par

Medíante algoritmos de IA se falsifican capturas de rendimientos “sin riesgo”, interfaces de software de selección cuantitativa de acciones o informes de análisis de inversión “de maestro”, y con el cebo de altos retornos se vende un servicio de inversión caro o se induce a la víctima a transferir fondos a cuentas controladas por los estafadores,\textbf{que al final desaparecen con el dinero}.

\booksubsection{Cómo desmontarlos — la “lupa” protectora}

La IA puede clonar voz y rostro, pero no es infalible; se pueden aprovechar varios fallos:

\begin{itemize}
  \item \textbf{Verificación por acción}: en una videollamada, pide al interlocutor que complete una secuencia de movimientos, como "pasa la mano rápidamente por delante de la cara", "gira la cabeza para mostrar el contorno de la oreja" o "parpadea 5 veces seguidas". Las técnicas de intercambio de rostro por IA actuales\textbf{suelen fallar}con oclusiones dinámicas complejas, movimientos rápidos y cambios de iluminación, produciendo
  \item \textbf{congelamientos, distorsiones faciales o bordes defectuosos}.\textbf{Verificación por contraseña privada}.
  \item \textbf{termina la llamada de inmediato y activa el protocolo de seguridad}.\textbf{Verificación secundaria}: sigue estrictamente el principio de "\textbf{cuatro no y una verificación}" de la Oficina de Ciberseguridad del Ministerío de Público: no creas en mensajes de desconocidos, no transfieras dinero a la ligera, no reveles información personal sensible, no difundas contenido sospechoso. Ante cualquier movimiento de dinero,
\end{itemize}

\booksubsection{devuelve la llamada desde un número ya conocido o confirma en persona}

\begin{itemize}
  \item \textbf{. Al transferir, elige el modo "ingreso en 24 horas" para disponer de tiempo suficiente para detectar el fraude y recuperar el dinero.}3. Timos clásicos (sigue habiendo que estar alerta):
  \item \textbf{Premíos / reembolso por tareas}: cebo; primero te dan un pequeño beneficio y luego te estafan una cantidad grande.
  \item \textbf{Estafa amorosa}: "matanza del cerdo"; tras cultivar el afecto, te estafan el dinero.
\end{itemize}

\booksubsection{Suplantación de un conocido}

\begin{itemize}
  \item : tras robar una cuenta, piden dinero prestado (en la era de la IA, aún más difícil de distinguir).
  \item 4. Protecciones integradas en el teléfono:
  \item Bloqueo de llamadas molestas.
\end{itemize}

\booksubsection{Protección contra robo de códigos de verificación.}

\textbf{Análisis de seguridad de los enlaces en SMS.}5. Fórmula antifraude:

\booksection{V. Resistirse al “doxxing” y a los ataques de ingeniería social: cortar la cadena de la fuga de privacidad}

Si te encuentras con un “asunto urgente”, mantén la calma 5 minutos; casi todos los fraudes delatan su cola en ese momento.\textbf{V. Resistirse al “doxxing” y a los ataques de ingeniería social: cortar la cadena de la fuga de privacidad}El principio central del “doxxing” (ataque de ingeniería social) es

\booksubsection{recopilar, analizar y recomponer sistemáticamente fragmentos digitales personales para construir un perfil completo}

Los métodos básicos de limpieza de EXIF se encuentran en el Capítulo 2 “Gestión de imágenes y sistema de archivos”. Antes de compartir fotos en redes sociales, es necesario además\textbf{1. Anonimización de fotos}Antes de compartir fotos en redes sociales, hay que usar herramientas de edición profesionales para\textbf{eliminar por completo los metadatos (EXIF)}, incluidos fecha de toma, modelo del dispositivo, coordenadas GPS, etc. La información clave que aparezca en la foto, como

\begin{itemize}
  \item números de puerta, matrículas, códigos de barras de billetes, números de identificación
  \item , debe difuminarse o pixelarse para no revelar involuntariamente tus movimientos.
  \item iOS: en [Configuración] > [Privacidad y seguridad] > [Cámara] puedes desactivar "Ubicación precisa"; al compartir, elige "Solo esta foto" para eliminar los metadatos.
\end{itemize}

\booksubsection{Android: algunas App de álbum (como Google Fotos) preguntan al compartir si se conserva la ubicación.}

Los métodos básicos de gestión de permisos se detallan en el Capítulo 4 “Gestión de permisos”. En escenarios de “prevención de doxxing”, es necesario además revisar periódicamente y desactivar los permisos de las apps de comercio electrónico y herramientas\textbf{2. Restricción de permisos de las App}Revisa periódicamente y desactiva los permisos no necesarios de micrófono, contactos y ubicación de las App de comercio y de herramientas, para evitar que “escuchen” conversaciones en segúndo plano o recopilen ilegalmente tu red de contactos; desinstala a tiempo las App que lleven mucho tiempo sin usarse pero que sigan teniendo permisos elevados, reduciendo desde la raíz los puntos de fuga de privacidad.

\begin{itemize}
  \item Cuando veas una solicitud de "permiso de contactos", pregúntate: ¿esta App realmente necesita conocer a todos mis contactos?
  \item Si llevas tres meses sin usar una App, desinstálala. Los permisos se restringen con el uso, no se aflojan.
\end{itemize}

\booksubsection{Seguridad de los dispositivos inteligentes}

Cámaras inteligentes, cerraduras inteligentes, altavoces inteligentes y otros dispositivos IoT deben\textbf{cambiar inmediatamente la contraseña por defecto}al usarse por primera vez; establece una contraseña independiente y robusta con mayúsculas, minúsculas, números y símbolos especiales. Las funciones de acceso remoto que no uses deben permanecer desactivadas; revisa periódicamente las actualizaciones de firmware y parchea las vulnerabilidades a tiempo,\textbf{para evitar que los hackers entren y espíen tu espacio privado}.

\begin{tcbwarn}{Advertencia con caso real}
En 2022, la policía de una localidad informó de que un hacker, aprovechando la contraseña por defecto de cámaras domésticas de cierta marca (muchos usuarios nunca la cambiaban), sustrajo imágenes de los dormitoríos de miles de familias y las vendió en internet.\textbf{No cambiar la contraseña por defecto = dejar la puerta de casa abierta}.

\end{tcbwarn}

\booksection{VI. Protección de la privacidad: tus datos son tu patrimonio}

Las fotos, el historial de chat y los movimientos de ubicación del teléfono son, en esencia, tus “activos digitales”. Si se filtran, en el mejor de los casos recibirás publicidad molesta; en el peor, podrás ser doxxeado, estafado o incluso chantajeado.

\booksubsection{Principio de “minimización” de permisos de App:}

\begin{itemize}
  \item \textbf{Ubicación}: "Permitir siempre" para transporte/mapas; "Permitir solo durante el uso" para el resto.
  \item \textbf{Contactos}: solo si es necesario (como añadir amigos en WeChat); en caso contrario, denegar.
  \item \textbf{Álbum / todas las fotos}: desde iOS 14+ se puede autorizar "Fotos seleccionadas"; Android 13+ también admite permisos detallados.
  \item \textbf{Micrófono / cámara}: desactivados por defecto, actívalos cuando los necesites. Si sospechas que te graban a escondidas, desliza el centro de control y comprueba si aparece un "punto verde/naranja".
  \item \textbf{Seguimiento en segúndo plano}: desde iOS 14.5+ se exige a las App que pregunten antes de usar IDFA; en Android se puede desactivar el "ID de publicidad".
\end{itemize}

\booksubsection{Gestión práctica de la privacidad:}

\begin{itemize}
  \item \textbf{Informe de privacidad (iOS)}: [Configuración] > [Privacidad y seguridad] > [Informe de privacidad de la App], para ver qué App acceden a tu ubicación, cámara y micrófono a escondidas.
  \item \textbf{Registro de uso de permisos (Android)}: [Configuración] > [Privacidad] > [Registro de uso de permisos], para ver quién invoca permisos sensibles y cuándo.
  \item \textbf{Limitar el seguimiento publicitarío}: en los ajustes del sistema, desactiva el "ID de publicidad" o "Limitar seguimiento publicitarío" para reducir la "publicidad precisa molesta".
\end{itemize}

\booksection{VII. Lista de “primeros auxilios” tras una pérdida o robo}

Si pierdes el teléfono, no te asustes; sigue este orden:

\begin{enumerate}
  \item \textbf{Congela la tarjeta SIM inmediatamente}: llama al operador para suspenderla y evitar que se use para recibir códigos de verificación.
  \item \textbf{Bloqueo remoto}: activa el "Modo perdido" a través del servicio en la nube (iCloud o la nube de la marca).
  \item \textbf{Cambia las contraseñas importantes}: cambia de inmediato la contraseña de todas las cuentas importantes que hayan iniciado sesión en el teléfono (pagos, banco, correo).
  \item \textbf{Denuncia}: presenta la denuncia en la comisaría y anota el número IMEI. Si recuperas el teléfono, también facilitará acreditar la titularidad.
  \item \textbf{Avisa a familiares y amigos}: evita que los estafadores pidan dinero prestado en tu nombre.
  \item \textbf{(Al final) Borrado remoto}: solo cuando confirmes que no hay forma de recuperarlo; tras el borrado, la localización deja de funcionar.
\end{enumerate}

\booksection{VIII. Seguridad de la tarjeta SIM y de la eSIM}

\begin{itemize}
  \item \textbf{Configura el PIN de la tarjeta SIM}: al activarlo, cada reinicio del teléfono o cambio de tarjeta exigirá el PIN, evitando que "saquen la tarjeta y la pongan en otro teléfono" para saltarse el bloqueo de pantalla.
  \item \textbf{Ventajas de la eSIM}: no hay tarjeta física que se pueda extraer, mayor seguridad; admite además el cambio entre varios números y es muy cómoda para viajar.
  \item \textbf{Riesgo de la "segúnda sustitución de tarjeta"}: cuidado con la estafa de "segúnda sustitución de tarjeta del operador": el estafador suplanta tu identidad para reponer la tarjeta SIM y hacerse con tu número. Configurar "la reposición de tarjeta requiere que el titular acuda en persona con su DNI a la oficina" es la defensa más eficaz.
\end{itemize}

\begin{tcbwarn}{Lo que nunca debes hacer}
\begin{itemize}
  \item Usar contraseñas débiles como cumpleaños / 123456 / qwerty;
  \item Enviar el código de verificación a nadie (ni siquiera al "servicio al cliente");
  \item Hacer clic en enlaces desconocidos de SMS, sobre todo de números extranjeros que empiezan por 00;
  \item Iniciar sesión en la banca online o en pagos bajo Wi-Fi público;
  \item Activar la "pantalla compartida" o el "control remoto" con desconocidos.
\end{itemize}

\end{tcbwarn}

\begin{tcbpractice}{Prueba a activar y simular el "bloqueo remoto"}
\begin{enumerate}
  \item Asegúrate de tener activada la función "Buscar mi teléfono".
  \item En el navegador de una computadora o en otro teléfono, inicia sesión en la cuenta del servicio en la nube correspondiente (iCloud.com o el sitio de la nube de la marca).
  \item Busca la función "Buscar dispositivo" y lanza una orden de "Reproducir sonido" o simula un "Bloqueo" sobre tu propio teléfono para familiarizarte con todo el proceso.
\end{enumerate}

\end{tcbpractice}
